\section{Open Problems and Discriminating Experiments}
\label{sec:open}

The largest uncertainties are actionable because each has a measurable descriptor. The proposals below condense the three experiment plans in the ideas directory; all precursor suggestions produced by the local two-stage model are explicitly marked as model predictions, pending experimental validation. They are research designs, not laboratory instructions.

\subsection{Lithium activity during phase formation}

High-temperature diffraction and multielement-loss studies show that lithium inventory, exposed thickness and processing path can change phase evolution and final composition \cite{gullbrekken2025phase,montoya2024lithium,moy2025irreversible}. The proposed test is a blocked $2\times2\times2$ interaction screen for Al--LLZO that varies lithium-addition order, compact thickness and cover state at a fixed nominal excess, followed by a smaller excess-lithium grid for the two strongest interactions. Every condition receives high-temperature or operando diffraction, density and a sister-specimen Li/Al/La/Zr assay; conductivity is compared only after density normalization. The decision rule is a reproducible cubic phase fraction, assay-corrected lithium retention and a conductivity change that survives furnace-block adjustment.

\paragraph{Model prediction, pending experimental validation.}
For Li$_{6.25}$Al$_{0.25}$La$_3$Zr$_2$O$_{12}$, the local precursor model ranks ZrO$_2$ + La$_2$O$_3$ + Al$_2$O$_3$ + Li$_2$CO$_3$ above alternative oxide or hydroxide sets. The ranking is a route suggestion only: hydroxide tokens require chemical identity and stoichiometry review before any experiment. The scientific test is the interaction contrast, not the ranking itself.

\subsection{Grain-boundary chemistry at matched density}

Co-doping and boundary-engineering studies suggest that a lower grain-boundary resistance can reflect a changed boundary composition rather than a more conductive garnet bulk \cite{ahmad2025reducing,golmohammad2024effects,du2024constructing,yao2026elimination}. A Ta--LLZO pilot should reproduce the Li$_2$CO$_3$ versus LiOH precursor comparison while selecting two furnace schedules that reach matched relative-density bands. Cross-sectional STEM--EDS/EELS, or a calibrated ToF--SIMS fallback, is required to quantify boundary composition or thickness; impedance and a standardized CCD test are then performed on the same density bands \cite{bae2026effect,klimpel2023standardizing}. The hypothesis is supported only if one calibrated boundary descriptor predicts both grain-boundary resistance and CCD after density and furnace block are included.

\paragraph{Model prediction, pending experimental validation.}
For Li$_{6.4}$La$_3$Zr$_{1.4}$Ta$_{0.6}$O$_{12}$, the local model ranks oxide precursors with LiOH or Li$_2$CO$_3$ as alternatives. The ambiguous LiHO token is not a validated reagent; the experiment must use a chemist-verified commercial precursor and record hydration state by thermal analysis.

\subsection{Surface layers after ultrafast sintering}

Depth-resolved spectroscopy identifies a Li$_2$O-rich surface layer in ultrafast-sintered LLZO, while thermal cleaning and membrane studies show that surface chemistry depends on atmosphere and post-heat treatment \cite{zhang2024unveiling,zhang2023on,zhang2023ultrafast,shi2025li2co3}. The proposed paired-sister design compares ultrafast and conventional sintering with 0, 1 and 2 h post-heating near the reported 900~$^{\circ}$C mitigation condition. One sister is reserved for as-sintered electrochemistry, one for depth chemistry and lithium assay, one for post-heat chemistry, and one for post-heat electrochemistry. XPS depth profiles, FIB--ToF--SIMS, carbonate spectroscopy, phase analysis and density are combined with the standardized CCD protocol \cite{klimpel2023standardizing}. Cleaning is considered beneficial only when the surface signal decreases, lithium retention and cubic phase are preserved, and the paired interface resistance or CCD improves.

\paragraph{Model prediction, pending experimental validation.}
The same local precursor ranking used for Al--LLZO favors oxide\slash carbonate mixtures, but it does not predict the thickness or chemistry of the post-sinter surface layer. Those quantities require direct depth calibration and an inert-transfer control.

These experiments are deliberately orthogonal: the first changes lithium activity, the second isolates boundary chemistry at matched density, and the third separates surface cleaning from bulk densification. Their measurements can therefore falsify a proposed mechanism instead of merely producing another optimized conductivity value.
